% from ega4-11-fr.tex
\oldpage[IV-3]{124}

\begin{enumerate}[label=\emph{(\roman*)}]
\item Soient $x$ un point de $X$, $x_\lambda$ sa projection canonique dans
$X_\lambda$. Pour que $\mathcal F$ soit $f$-plat au point $x$, il faut et il
suffit qu'il existe $\lambda\geq\alpha$ tel que $\mathcal F_\lambda$ soit
$f_\lambda$-plat au point $x_\lambda$.
\item Pour que $\mathcal F$ soit $f$-plat, il faut et il suffit qu'il existe
$\lambda\geq\alpha$ tel que $\mathcal F_\lambda$ soit $f_\lambda$-plat.
\end{enumerate}
\end{theorem}

On peut supposer que $S_\alpha=S_0$~; comme $Y_0$ est de présentation finie
sur $S_0$, $Y_0$ est quasi-compact, et $f_0:X_0\to Y_0$ est un morphisme de
présentation finie (1.6.2, (v)), donc on peut aussi se borner au cas où
$S_0=Y_0$. En outre, en vertu de la quasi-compacité de $S_0$ et du fait que
l'ensemble d'indices $L$ est filtrant, on peut se borner au cas où
$S_0=\operatorname{Spec}(A_0)$ est affine. De plus $X_0$ est quasi-compact,
donc le même raisonnement montre qu'on peut aussi supposer
$X_0=\operatorname{Spec}(B_0)$ affine~; on a alors
$\mathcal F_0=\widetilde M_0$, où $M_0$ est un $B_0$-module de présentation
finie, et l'énoncé (11.2.6) est dans ce cas équivalent au suivant (compte tenu
de ($\mathbf 0_{\mathrm I}$, 6.3.1))~:

\begin{corollary}[11.2.6.1]
\label{IV.11.2.6.1-fr}
Soient $A_0$ un anneau, $(A_\lambda)_{\lambda\in L}$ un système inductif
filtrant de $A_0$-algèbres, $B_0$ une $A_0$-algèbre de présentation finie,
$M_0$ un $B_0$-module de présentation finie. On pose
\[
\begin{gathered}
B_\lambda=B_0\otimes_{A_0}A_\lambda,\qquad
M_\lambda=M_0\otimes_{A_0}A_\lambda=M_0\otimes_{B_0}B_\lambda,\\
A=\varinjlim A_\lambda,\qquad
B=\varinjlim B_\lambda=B_0\otimes_{A_0}A,\\
M=\varinjlim M_\lambda=M_0\otimes_{A_0}A=M_0\otimes_{B_0}B.
\end{gathered}
\]
\begin{enumerate}[label=\emph{(\roman*)}]
\item Soit $\mathfrak p$ un idéal premier de $B$, et pour tout $\lambda$,
soit $\mathfrak p_\lambda$ son image réciproque dans $B_\lambda$. Pour que
$M_{\mathfrak p}$ soit un $A$-module plat, il faut et il suffit qu'il existe
$\lambda$ tel que $(M_\lambda)_{\mathfrak p_\lambda}$ soit un
$A_\lambda$-module plat.
\item Pour que $M$ soit un $A$-module plat, il faut et il suffit qu'il existe
$\lambda$ tel que $M_\lambda$ soit un $A_\lambda$-module plat.
\end{enumerate}
\end{corollary}

On doit seulement démontrer que les conditions sont nécessaires (2.1.4).
Nous procéderons en plusieurs étapes.

\medskip
\noindent\textbf{I)} \emph{Réduction au cas où les $A_\lambda$ sont
noethériens.} En vertu de (8.9.1), il existe un sous-anneau $A'_0$ de $A_0$,
qui est une $\mathbf Z$-algèbre de type fini, une $A'_0$-algèbre de type fini
$B'_0$ et un $B'_0$-module de type fini $M'_0$ tels que l'on ait
$B_0=B'_0\otimes_{A'_0}A_0$ et $M_0=M'_0\otimes_{A'_0}A_0$~; comme on a
$B_\lambda=B'_0\otimes_{A'_0}A_\lambda$ et
$M_\lambda=M'_0\otimes_{A'_0}A_\lambda$, on peut remplacer $A_0$, $B_0$,
$M_0$ par $A'_0$, $B'_0$, $M'_0$ dans l'énoncé de (11.2.6.1), en considérant
les $A_\lambda$ comme des $A'_0$-algèbres~; on peut donc déjà supposer que
$A_0$ est noethérien. Soit $H$ l'ensemble des couples
$(\lambda,C_\lambda)$, où $C_\lambda$ est une sous-$A_0$-algèbre de
\emph{type fini} de $A_\lambda$~; ordonnons $H$ en posant
$(\lambda,C_\lambda)\leq(\mu,D_\mu)$ si $\lambda\leq\mu$ et si
l'homomorphisme $\varphi_{\mu\lambda}:A_\lambda\to A_\mu$ est tel que
$\varphi_{\mu\lambda}(C_\lambda)\subset D_\mu$~; pour cet ordre $H$ est
filtrant croissant, car si $(\lambda,C_\lambda)$ et $(\mu,D_\mu)$ sont deux
éléments quelconques de $H$, on les majorera par un $(\nu,E_\nu)$ en prenant
$\nu\geq\lambda$, $\nu\geq\mu$ dans $L$, puis $E_\nu$ égal à la
sous-$A_0$-algèbre de $A_\nu$ engendrée par
$\varphi_{\nu\lambda}(C_\lambda)$ et $\varphi_{\nu\mu}(D_\mu)$. Pour un
élément $\xi=(\lambda,C_\lambda)$ de $H$, on posera $A_\xi=C_\lambda$, et
pour $\xi\leq\eta=(\mu,D_\mu)$ (donc $\lambda\leq\mu$ et
$\varphi_{\mu\lambda}(C_\lambda)\subset D_\mu$),
$\varphi_{\eta\xi}:A_\xi\to A_\eta$ sera la restriction à $C_\lambda$ de
$\varphi_{\mu\lambda}$, considéré comme homomorphisme dans $D_\mu$~; il est
clair que l'on obtient ainsi un système inductif filtrant de $A_0$-algèbres.
On pose $B_\xi=B_0\otimes_{A_0}A_\xi$,
$M_\xi=M_0\otimes_{A_0}A_\xi$~; cette fois les $A_\xi$ sont
noethériens~; en outre la formule de la double limite inductive (Bourbaki,
\emph{Alg.}, chap.~II, 3$^e$ éd., \S~6, n$^\circ$~4, prop.~7) prouve que
l'on a encore
$\varinjlim_HA_\xi=A$, $\varinjlim_HB_\xi=B$,
$\varinjlim_HM_\xi=M$. Supposons (11.2.6.1) prouvé pour le système inductif
$(A_\xi)_{\xi\in H}$~; soit $\mathfrak p$ un idéal premier de $B$,

